\chapter*{서문: 왜 지금 AI 거버넌스인가?}
\addcontentsline{toc}{chapter}{서문: 왜 지금 AI 거버넌스인가?}
\markboth{서문}{서문}

\section*{현장에서 목격한 AI의 '빛과 그림자'}

필자는 수년간 다양한 산업군의 데이터 패브릭 및 AI 프로젝트를 관리하며 수많은 기업들의 디지털 전환 과정을 목격해 왔다. 그 과정에서 얻은 가장 큰 깨달음은, AI의 성능보다 더 중요한 것은 \textbf{'그 기술을 담는 그릇의 견고함'}이라는 사실이다.

현장에서 본 많은 기업들은 AI 도입 초기, 화려한 정확도 수치와 알고리즘의 최신성에 열광한다. 하지만 프로젝트가 실제 운영 단계에 접어들면 상황은 급변한다. 
\begin{itemize}
    \item "학습 데이터에 편향이 섞여 특정 고객층에게 불리한 결과가 나온다면 누가 책임질 것인가?"
    \item "모델의 결정 근거를 설명하지 못해 규제 당국이나 고객의 항의에 대응하지 못하면 어떻게 할 것인가?"
    \item "운영 중인 모델의 성능이 소리 없이 저하(드리프트)되어 비즈니스에 큰 손실을 입히고 있지는 않은가?"
\end{itemize}

이러한 질문들은 기술적인 문제를 넘어 조직의 생존과 직결되는 전략적 리스크이다. 필자는 공공부문 데이터 거버넌스 체계 수립과 대형 금융사의 AI 플랫폼 구축 프로젝트를 리딩하며, 이러한 리스크를 제안 단계부터 선제적으로 관리하지 못했을 때 발생하는 엄청난 기회비용과 사회적 불신을 뼈아프게 경험했다.

\section*{실무자를 위한 '해답지'를 만들다}

오늘날 AI 거버넌스에 관한 논의는 활발하지만, 정작 실무 현장에서는 "그래서 구체적으로 무엇부터 해야 하는가?"라는 질문에 답해주는 가이드가 부족하다. 해외의 추상적인 윤리 원칙이나 복잡한 법조문 나열만으로는 현업의 병목 현상을 해결할 수 없다.

본 저서는 이러한 문제의식에서 출발했다. 단순히 '무엇(What)'을 해야 하는지를 나열하는 대신, \textbf{'어떻게(How)'} 구현하고 운영할 것인지에 초점을 맞추었다. 필자가 실무를 수행하며 다듬어온 데이터 패브릭 기반의 거버넌스 아키텍처, MLOps와 거버넌스의 통합 방법론, 그리고 실제 알고리즘 감사 과정에서 사용했던 체크리스트들을 가감 없이 담았다.

특히 2026년부터 본격적으로 시행된 \textbf{'AI 기본법'}의 요구사항을 실무적으로 어떻게 해석하고 대응해야 하는지, 한국 기업 환경에 최적화된 로드맵은 무엇인지에 대해 심도 있게 다루었다.

\section*{거버넌스는 규제가 아니라 성장 엔진이다}

이 책을 쓰면서 필자가 가장 강하게 전달하고 싶었던 메시지는 하나다. \textbf{AI 거버넌스는 AI의 발목을 잡는 규제가 아니라, 더 빠르게 달리게 해주는 에이블러(Enabler)다.}

이것은 낙관적 수사가 아니라 데이터가 뒷받침하는 사실이다. IBM 기업가치연구소(2026)가 전 세계 기업을 분석한 결과, \textbf{AI 거버넌스를 완전히 내재화한 기업은 그렇지 않은 기업보다 매출 성장을 보고할 확률이 약 4배 높다}(58\% vs 15\%). PwC의 2026년 AI 전망 보고서도 같은 결론을 내린다: ``거버넌스를 먼저 구축하고, 인력을 준비시키고, 효과 없는 것을 멈출 규율을 갖춘 조직들이 모든 지표에서 경쟁자를 앞서고 있다.''\cite{pwc2026aivalue}

왜 그럴까? 거버넌스가 있는 조직은 \textbf{``Yes''를 더 빠르게 말할 수 있기 때문이다.} 통제 구조가 있으면 리더는 야심 찬 AI 프로젝트에 더 빨리 승인을 내릴 수 있다. 위험 관리 체계가 작동한다는 신뢰가 있을 때, 혁신은 더 빠르게, 더 멀리 나아간다. 반대로 거버넌스 없는 조직에서 AI 도입의 진짜 장벽은 기술이 아니라 신뢰의 부재다.

이 책에서 소개하는 사례들---네덜란드의 SyRI, 호주의 Robodebt, 그리고 국내 공공·금융·의료 프로젝트의 실패들---은 모두 같은 이야기를 한다. \textbf{거버넌스의 부재가 기술의 실패보다 더 치명적이다.} 그러나 반대 방향의 이야기도 동등하게 중요하다. 거버넌스를 조기에 내재화한 기업들은 규제 리스크를 피했을 뿐 아니라, 신뢰 자산을 바탕으로 더 빠른 시장 진입과 더 높은 고객 충성도를 확보했다.

\section*{이 책을 읽는 독자들께}

AI 거버넌스는 준수해야 할 체크리스트가 아니다. 조직이 AI로 지속가능한 경쟁 우위를 만들어가는 \textbf{전략적 역량}이다.

이 책이 AI 도입을 고민하는 경영진에게는 ``왜 거버넌스가 투자 대비 수익이 가장 높은 AI 전략인가''에 대한 확신을, 실무자들에게는 바로 내일 현장에서 꺼내 쓸 수 있는 \textbf{거버넌스 툴킷}을 제공하기를 바란다. 그리고 대한민국 AI 생태계가, 규제를 피하는 방향이 아니라 신뢰를 축적하는 방향으로 성장하는 데 이 책이 작은 이정표가 되기를 희망한다.

\vspace{2cm}
\begin{flushright}
\textbf{저자 씀} \\[0.3em]
곽두일 \\
AI 컨설턴트
\end{flushright}
